\chapter{Niezmienniki Seiberga--Wittena i egzotyczność rozmaitości Parka}
\label{ch:tg-park-sw}

W poprzednim rozdziale opisaliśmy konstrukcję gładkiej,
zamkniętej, orientowanej i po prostu spójnej
czterorozmaitości $X_7$, korzystając z geometrycznych
wyników Parka i Cassona--Harera.
Wiemy, że ma ona formę przecięcia izometryczną z formą
$\langle1\rangle\oplus7\langle-1\rangle$, a zatem jest
homeomorficzna z $\mathbb{CP}^2\#7\overline{\mathbb{CP}}^{\,2}$.
Nie wynika stąd jednak, że obie rozmaitości są dyfeomorficzne.
Pokażemy, dlaczego nie są, korzystając z teorii Seiberga--Wittena.

Wymiar dodatniej części formy przecięcia wynosi tu $b_2^+=1$.
Przy tym założeniu liczba otrzymana przez zliczenie rozwiązań
równań Seiberga--Wittena zależy na ogół od \emph{komory}.
Z tego powodu musimy szczególnie starannie oddzielić niezmiennik
z małą perturbacją, który dla rozważanej klasy charakterystycznej
jest niezależny od metryki, od niezmiennika w \emph{komorze
symplektycznej}. Wybór właściwej komory jest zasadniczym krokiem
rozróżnienia struktur gładkich.

\section{Klasy sfer i ich przecięcia: uzupełnienie konstrukcji}
\label{sec:tg-park14-spheres}

Przypomnijmy bazę $h,e_1,\ldots,e_{13}$ grupy
$H_2(R;\mathbb Z)$, gdzie $R=E(1)\#4\overline{\mathbb{CP}}^{\,2}$,
$h^2=1$, $e_i^2=-1$ i wszystkie pozostałe przecięcia między
różnymi elementami bazy są zerowe. Oznaczmy
$f=3h-e_1-\cdots-e_9$. Wybierzmy pięć klas
\[
\begin{aligned}
 v_1&=e_4-e_7, & v_2&=e_1-e_4, \\
 v_3&=h-e_1-e_2-e_3, & v_4&=e_2-e_5, \\
 v_5&=e_5-e_9,
\end{aligned}
\]
oraz $s=4f+e_9-2(e_{10}+e_{11}+e_{12}+e_{13})$.
Park wykorzystuje realizację klas $v_i$ przez pięć sfer
jednego włókna osobliwego typu $\widetilde E_6$, a klasy $s$
przez szóstą sferę uzyskaną po rozwiązaniu przecięć.
Geometryczne istnienie takiej realizacji przyjmujemy z jego pracy.

\begin{stwierdzenie}[Ustalenie numeracji]
\label{prop:tg-park14-classes}
Połóżmy $u_0:=s$ oraz $u_j:=v_{6-j}$ dla $1\le j\le5$.
Wówczas $u_0^2=-9$, $u_j^2=-2$ dla $j\ge1$,
$u_j\cdot u_{j+1}=1$ dla $0\le j\le4$, a pozostałe
przecięcia różnych sfer są zerowe.
\end{stwierdzenie}

\begin{proof}
Jest to dokładnie wynik lematu
\ref{lem:tg-park-S-classes} z poprzedniego rozdziału,
po odwróceniu numeracji pięciu sfer o samoprzecięciu $-2$.
Podane tam obliczenie wykorzystuje jedynie diagonalną
formę przecięcia na $H_2(R;\mathbb Z)$.
\end{proof}

\begin{uwaga}[Zakres dowodu]
Stwierdzenie jest \emph{rachunkiem homologicznym}.
Nie dowodzi, że dowolne sześć klas o takiej macierzy
ma reprezentantów będących zanurzonymi sferami
symplektycznymi o zadanym układzie przecięć.
To właśnie niezależna geometryczna treść konstrukcji Parka.
Jest ona potrzebna do zastosowania twierdzenia Symington.
\end{uwaga}

\section{Klasa kanoniczna i odwrotność macierzy plumbingowej}
\label{sec:tg-park14-canonical}

Niech $(R,\omega)$ będzie symplektyczną powierzchnią
racjonalną z konfiguracją Parka, a $K_R=-c_1(TR,J)$
jej klasą kanoniczną dla zgodnej struktury prawie
zespolonej $J$. Utożsamiamy klasy homologii i kohomologii
przez dualność Poincarégo. Dla standardowej klasy
kanonicznej, użytej w konstrukcji Parka, wzór na
klasę kanoniczną po rozdmuchaniu ($K_{\widetilde X}
=\pi^*K_X+\operatorname{PD}[E]$ dla krzywej
wyjątkowej $E$) daje
\[
K_R=-3h+e_1+\cdots+e_{13}
    =-f+e_{10}+e_{11}+e_{12}+e_{13}.
\]
Istotnie, $K_{\mathbb{CP}^2}=-3h$ (z ciągu
Eulera dla wiązki stycznej przestrzeni rzutowej),
a trzynaście rozdmuchań dodaje w bazie całkowitych
transformat po jednej klasie wyjątkowej.
W szczególności $K_R\cdot h=-3$ i
$K_R\cdot e_j=-1$. Stąd
$K_R\cdot(e_j-e_k)=0$ oraz
$K_R\cdot(h-e_1-e_2-e_3)=-3+3=0$,
czyli $K_R\cdot u_i=0$ dla $1\le i\le5$.
Aby policzyć $K_R\cdot u_0$, można wykorzystać
formułę adjunkcji dla symplektycznej sfery $u_0$:
$-2=u_0^2+K_R\cdot u_0=-9+K_R\cdot u_0$.
Dostajemy $K_R\cdot u_0=7$. Ten sam wynik otrzymujemy bezpośrednio:
$f\cdot s=1$, a $e_j\cdot s=2$ dla $10\le j\le13$,
więc $K_R\cdot s=-1+4\cdot2=7$.

\begin{lemat}[Macierz odwrotna konfiguracji $C_7$]
\label{lem:tg-park14-inverse}
Niech $Q_7$ będzie macierzą przecięcia konfiguracji $C_7$.
Jej odwrotność jest określona nad $\mathbb Q$ i ma pierwszy
wiersz postaci
\[
(Q_7^{-1})_{0*}=-\frac1{49}(6,5,4,3,2,1).
\]
\end{lemat}

\begin{proof}
Jest to pierwszy wiersz pełnej odwrotności z lematu
\ref{lem:tg-park-Q-inverse-full}. Dla niezależnej kontroli
przypomnijmy, że $\det Q_7=49$, więc macierz jest
odwracalna nad $\mathbb Q$. Aby znaleźć pierwszy wiersz,
szukamy wektora $z=(z_0,\ldots,z_5)^T$ spełniającego
$Q_7z=(1,0,\ldots,0)^T$; symetria $Q_7$ kończy dowód.
Pięć ostatnich równań to kolejno
$z_{j-1}-2z_j+z_{j+1}=0$ dla $1\le j\le4$
oraz $z_4-2z_5=0$. Wynika z nich
$z_j=(6-j)z_5$ dla $0\le j\le5$.
Pierwsze równanie $-9z_0+z_1=1$ daje
$-49z_5=1$. Otrzymujemy
$z_j=-(6-j)/49$, jak w tezie.
\end{proof}

\begin{stwierdzenie}[Rachunek parowania po wycięciu konfiguracji]
\label{prop:tg-park14-splitting}
Niech $V=R\setminus\operatorname{int}C_7$,
a $X_7=V\cup B_7$ będzie symplektyczną wymierną
kontrakcją, przy której formy symplektyczne na $V$
są identyfikowane przez symplektomorfizm.
Niech $\alpha_i:=[\omega]\cdot u_i$ dla $0\le i\le5$.
Wtedy
\[
K_{X_7}\cdot[\omega_7]
 =K_R\cdot[\omega]+
   \frac17(6\alpha_0+5\alpha_1+4\alpha_2
                    +3\alpha_3+2\alpha_4+\alpha_5).
\]
\end{stwierdzenie}

\begin{proof}
Ponieważ $H_1(\partial C_7;\mathbb Q)
=H_2(\partial C_7;\mathbb Q)=0$,
ciąg Mayera--Vietorisa daje, nad $\mathbb R$,
izomorfizmy przez ograniczenie:
\[
H^2(R;\mathbb R)\cong H^2(V;\mathbb R)\oplus H^2(C_7;\mathbb R),
\qquad H^2(X_7;\mathbb R)\cong H^2(V;\mathbb R).
\]
Pracujemy nad $\mathbb R$, ponieważ klasa formy
symplektycznej nie musi być wymierna. Na każdym
składniku $W$ o brzegu będącym wymierną sferą
homologiczną mapa $H^2(W,\partial W;\mathbb R)
\to H^2(W;\mathbb R)$ jest izomorfizmem.
Parowanie na $H^2(W;\mathbb R)$ definiujemy
przez podniesienie do kohomologii względnej
i ewaluację iloczynu na klasie fundamentalnej
względnej. Reprezentanci względni mogą być
wybrani z nośnikiem we wnętrzu; przez przedłużenie
zerem uzyskujemy ortogonalny rozkład parowania
w powyższych izomorfizmach.
Oznaczmy przez $\gamma_i$ bazę $H^2(C_7;\mathbb R)$
sprzężoną z bazą $u_i$, czyli
$\langle\gamma_i,u_j\rangle=\delta_{ij}$. Ograniczenie klasy kanonicznej ma
wektor współrzędnych $(7,0,\ldots,0)$, a $[\omega]$
ma współrzędne $(\alpha_0,\ldots,\alpha_5)$.
Formę przecięcia w bazie sprzężonej reprezentuje
$Q_7^{-1}$: obrazy klas sfer pod dualnością względną
mają w bazie $\gamma_i$ kolumny $Q_7$, więc
przekształcenie macierzy parowania z bazy tych
obrazów do bazy $\gamma_i$ daje
$(Q_7^{-1})^TQ_7Q_7^{-1}=Q_7^{-1}$.
Lemat \ref{lem:tg-park14-inverse} daje
\[
K_R|_{C_7}\cdot[\omega]|_{C_7}
  =7(1,0,\ldots,0)Q_7^{-1}\alpha
  =-\frac17\sum_{j=0}^{5}(6-j)\alpha_j.
\]
Przy wymiernej kontrakcji $H^2(B_7;\mathbb R)=0$.
Symplektomorfizm dopełnień identyfikuje ograniczenia
form symplektycznych. Identyfikuje także klasy kanoniczne:
przenosi zgodną strukturę prawie zespoloną na zgodną
strukturę prawie zespoloną, a przestrzeń takich struktur
jest ściągalna. Naturalność pierwszej klasy Cherna daje
więc $K_R|_V=K_{X_7}|_V$ w tym utożsamieniu.
Jeżeli konstrukcja wymaga uprzedniego dostosowania
otoczenia sfer, przez $\omega$ rozumiemy formę po tym
dostosowaniu; zakładamy identyfikację dopełnień jak
w tezie stwierdzenia. Kula $B_7$ nie wnosi wkładu
do parowania nad $\mathbb R$. Z ortogonalnego rozkładu
otrzymujemy więc
$K_{X_7}\cdot[\omega_7]
=K_R\cdot[\omega]-K_R|_{C_7}\cdot[\omega]|_{C_7}$.
Podstawiając obliczoną wartość, otrzymujemy tezę.
\end{proof}

\begin{uwaga}[Kontrola indeksu klasy kanonicznej]
W konfiguracji Parka $v=(7,0,\ldots,0)^T$ dla
$K_R$ oraz $-v$ dla klasy wyznacznikowej $-K_R$
kanonicznej struktury $\mathrm{Spin}^c$.
Z lematu \ref{lem:tg-park14-inverse} otrzymujemy
$v^TQ_7^{-1}v=49(-6/49)=-6$.
Jest to dokładnie warunek zachowania wymiaru
ze stwierdzenia \ref{prop:tg-sw-general-index-pairing}
dla łańcucha długości sześć.
Przed kontrakcją $K_R^2=9-13=-4$,
a po kontrakcji $K_{X_7}^2=-4+6=2$,
zgodnie z twierdzeniem \ref{thm:tg-park-numbers}
i wzorem dla struktury prawie zespolonej.
Nie daje to samo w sobie niezerowości niezmiennika:
dla $b_2^+=1$ nadal trzeba ustalić właściwą komorę.
\end{uwaga}

\section{Ścisła dodatniość parowania \texorpdfstring{$K_{X_7}\cdot[\omega_7]$}{K omega}}
\label{sec:tg-park14-positivity}

Kluczowa pozostaje możliwość wyboru struktury symplektycznej,
zgodnej z opisaną konfiguracją, dla której współrzędne klasy
symplektycznej spełniają nierówności potrzebne w rachunku.
Wykorzystamy tu następującą postać twierdzeń o stożku
symplektycznym powierzchni racjonalnych, używaną przez Parka.

\begin{twierdzenie}[Geometryczne założenie o klasie symplektycznej; wynik zewnętrzny]
\label{thm:tg-park14-cone}
W konstrukcji Parka istnieje zgodna struktura symplektyczna
na $R$, dla której zanurzone sfery konfiguracji $C_7$
są symplektyczne, a klasa $[\omega]\in H^2(R;\mathbb R)$
ma reprezentację
\[
[\omega]=a h-\sum_{j=1}^{13}b_j e_j,
\qquad
 a\ge b_1\ge\cdots\ge b_{13}\ge0,
\qquad 3a>\sum_{j=1}^{13}b_j.
\]
Przyjmujemy standardową klasę kanoniczną opisaną
w poprzedniej sekcji. Istnienie jednocześnie takiej formy
i takiej symplektycznej realizacji konfiguracji przyjmujemy
jako dane geometryczne użyte w dowodzie twierdzenia 3.1
Parka (arXiv:math/0311395v3, s.~8--9); lemat 3.2
na s.~7 dotyczy wyboru klasy z teorii stożka
symplektycznego. Sam lemat 3.2 nie dowodzi
zgodności z ustalonymi sferami. Używamy nierówności
nieostrych $b_j\ge0$, dokładnie jak w pracy Parka.
\end{twierdzenie}

\begin{uwaga}[Zakres założenia geometrycznego]
Powyższa teza dotyczy konkretnej realizacji sfer w $R$.
Nie wolno wybrać dowolnych dodatnich współczynników
$a,b_j$ i na podstawie samych nierówności twierdzić,
że reprezentują one formę symplektyczną, względem której
wszystkie sfery konfiguracji są symplektyczne.
Istnienie formy, jej zgodność z ustalonymi klasami sfer
i znaki ich pól symplektycznych muszą zostać zweryfikowane
w rzeczywistej konstrukcji geometrycznej.
\end{uwaga}

\begin{przyklad}[Nierówności stożkowe nie zastępują geometrii sfer]
\label{ex:tg-park14-cone-not-sufficient}
Rozważmy formalną klasę
$[\beta]=a h-\sum_{j=1}^{13}b_j e_j$,
gdzie $a=10$, $b_1=b_2=b_3=4$
oraz $b_4=\cdots=b_{13}=1$.
Wtedy
\[
a\ge b_1\ge\cdots\ge b_{13}>0,
\qquad 3a=30>22=\sum_{j=1}^{13}b_j.
\]
Spełnione są zatem wszystkie wymienione
wyżej nierówności współczynników.
Mimo to dla klasy sfery
$v_3=h-e_1-e_2-e_3$
otrzymujemy
\[
[\beta]\cdot v_3=a-b_1-b_2-b_3=10-12=-2<0.
\]
Gdyby $v_3$ była symplektyczną sferą
z orientacją ustaloną przez formę
reprezentującą $[\beta]$, miałaby
ściśle dodatnie pole
$\int_{v_3}\beta>0$. Obliczenie
wyklucza więc taką interpretację
tej formalnej klasy.

Co więcej, ta formalna klasa ma dodatni kwadrat:
$[\beta]^2=100-3\cdot16-10=42>0$.
Nawet dołączenie warunku dodatniego kwadratu nie
zapewnia zatem dodatniości pól ustalonych sfer.

Przykład nie przeczy twierdzeniu
\ref{thm:tg-park14-cone}, które
\emph{zakłada jednocześnie}
istnienie odpowiedniej formy
i zgodnej z nią konfiguracji.
Pokazuje natomiast, dlaczego nie
wolno zastępować geometrycznego
warunku dodatniości pól sfer
samymi nierównościami
$a\ge b_1\ge\cdots\ge b_{13}\ge0$
oraz $3a>\sum b_j$.
\end{przyklad}

Dla takiego wyboru wyprowadzimy znak parowania
$K_{X_7}\cdot[\omega_7]$ przez całkowicie jawny rachunek.

\begin{twierdzenie}[Nierówność kluczowa w dowodzie Parka]
\label{thm:tg-park14-positive}
Niech $X_7$ pochodzi z symplektycznej wymiernej kontrakcji
w poprzednim rozdziale i niech $\omega_7$ odpowiada klasie
$[\omega]$ z twierdzenia \ref{thm:tg-park14-cone}.
Wtedy $K_{X_7}\cdot[\omega_7]>0$.
\end{twierdzenie}

\begin{proof}
Wprowadźmy skrót $B:=b_{10}+b_{11}+b_{12}+b_{13}$.
Z klas $u_0=s$, $u_1=v_5$, $u_2=v_4$, $u_3=v_3$,
$u_4=v_2$ i $u_5=v_1$ otrzymujemy bezpośrednio
\[
\begin{aligned}
\alpha_0&=12a-4\sum_{j=1}^{9}b_j+b_9-2B,
 &\alpha_1&=b_5-b_9,\\
\alpha_2&=b_2-b_5,
 &\alpha_3&=a-b_1-b_2-b_3,\\
\alpha_4&=b_1-b_4,
 &\alpha_5&=b_4-b_7.
\end{aligned}
\]
Na przykład w wyrażeniu dla $\alpha_0$ współczynnik
przy $b_9$ to $-3$, a nie $-4$, ponieważ klasa
$s=4f+e_9-2\sum_{j=10}^{13}e_j$ zawiera dodatkowy
składnik $+e_9$. Ponadto
$K_R\cdot[\omega]=-3a+\sum_{j=1}^{13}b_j$.

Korzystając ze stwierdzenia \ref{prop:tg-park14-splitting}
i grupując wyrazy, otrzymujemy
\[
\begin{aligned}
7K_{X_7}\cdot[\omega_7]
={}&54a-18b_1-16b_2-20b_3-18b_4-16b_5\\
 &-17b_6-18b_7-17b_8-16b_9-5B.
\end{aligned}
\]
Nie jest jeszcze oczywiste, że prawa strona jest dodatnia.
Rozbijmy ją zatem na dwie części:
\[
\begin{aligned}
7K_{X_7}\cdot[\omega_7]
={}&18\left(3a-\sum_{j=1}^{13}b_j\right)\\
 &+2(b_2-b_3)+2b_5+b_6+b_8+2b_9+13B.
\end{aligned}
\]
Pierwszy nawias jest ściśle dodatni z twierdzenia
o wyborze klasy symplektycznej. Pozostałe wyrazy
są nieujemne, ponieważ $b_2\ge b_3$ oraz $b_j\ge0$.
W szczególności nie trzeba zakładać ścisłej
dodatniości każdego $b_j$ ani ścisłych nierówności
w całym uporządkowanym ciągu współczynników.
Warunki te służą wyłącznie dowodowi algebraicznemu;
dodatniość pól wszystkich sfer wynika oddzielnie
z ich symplektycznej realizacji.
Wynika stąd $K_{X_7}\cdot[\omega_7]>0$.
\end{proof}

\begin{wniosek}[Ilościowe oszacowanie parowania kanonicznego]
\label{cor:tg-park14-quantitative-positive}
Niech $a,b_1,\ldots,b_{13}$ będą liczbami
rzeczywistymi spełniającymi
$b_2\ge b_3$, $b_j\ge0$ dla wszystkich
$1\le j\le13$ oraz
$3a\ge\sum_{j=1}^{13}b_j$.
Zakładamy ponadto $B:=b_{10}+b_{11}+b_{12}+b_{13}>0$.
Niech wyrażenie $P$ będzie określone
wzorem algebraicznym na
$K_{X_7}\cdot[\omega_7]$
z dowodu Twierdzenia~\ref{thm:tg-park14-positive}.
Wówczas
\[
P\ge\frac{13}{7}B>0.
\]
W szczególności do samego
\emph{algebraicznego} dowodu dodatniości
wystarcza słabsze założenie
$3a\ge\sum b_j$.
Nie twierdzimy, że każda lista liczb
spełniająca te nierówności jest
realizowana przez zgodną formę
symplektyczną na $R$.
\end{wniosek}

\begin{proof}
Obliczenie z poprzedniego dowodu daje
\[
7P=18\Bigl(3a-\sum_{j=1}^{13}b_j\Bigr)
+2(b_2-b_3)+2b_5+b_6+b_8+2b_9+13B.
\]
Na mocy założeń wszystkie składniki
przed $13B$ są nieujemne, a $B>0$.
Otrzymujemy $7P\ge13B>0$,
czyli żądane oszacowanie.
Jest to wyłącznie konsekwencja
rachunku liniowego; własności
geometryczne \emph{wymagane do
interpretacji} $P$ pozostają
wynikiem zewnętrznym.
\end{proof}

\begin{przyklad}[Kontrola algebraiczna wzoru]
Przyjmijmy wyłącznie do sprawdzenia rachunku
$a=40$ oraz $(b_1,\ldots,b_{13})=(13,12,11,10,9,8,7,6,5,1,1,1,1)$.
Wtedy $3a-\sum b_j=35>0$, a wektor przecięć
$(\alpha_0,\ldots,\alpha_5)=(153,4,3,4,3,3)$.
Zatem $K_R\cdot[\omega]=-35$ oraz
\[
K_{X_7}\cdot[\omega_7]
 =-35+\frac{6\cdot153+5\cdot4+4\cdot3
                    +3\cdot4+2\cdot3+3}{7}
 =\frac{726}{7}>0.
\]
Można także sprawdzić wynik przez bezpośrednie
rzutowanie ortogonalne: jeśli $U$ ma kolumny $u_i$
w bazie $h,e_1,\ldots,e_{13}$, a
$J_0=\operatorname{diag}(1,-1,\ldots,-1)$,
to $U^TJ_0U=Q_7$. Operator
$P_V=I-UQ_7^{-1}U^TJ_0$ jest rzutem na
dopełnienie ortogonalne kolumn $U$:
$P_VU=0$, $U^TJ_0P_V=0$ i
$P_Vx=x$ dla $U^TJ_0x=0$.
Zatem dla wektorów współrzędnych $k,w$ klas
$K_R,[\omega]$ parowanie po kontrakcji wynosi
$k^TJ_0P_Vw$. Dla samej klasy kanonicznej
dostajemy niezależną kontrolę jej kwadratu:
\[
K_{X_7}^2=K_R^2-(7,0,\ldots,0)Q_7^{-1}
(7,0,\ldots,0)^T=-4+6=2.
\]
Liczby te sprawdzają tożsamość algebraiczną;
nie twierdzimy, że dowolna zadana lista współczynników
opisuje formę symplektyczną zachowującą cały geometryczny
układ sfer Parka.
\end{przyklad}

\begin{center}
\begin{tikzpicture}[>=Stealth,scale=.95,line cap=round,line join=round]
  \node[draw=tgBlue,fill=tgBlue!8,rounded corners=4pt,
     minimum width=3.1cm,minimum height=1.1cm,align=center] (r)
     at (-4.0,0) {\small $R=E(1)\#4\overline{\mathbb{CP}}^{\,2}$\\[-1pt]
             \scriptsize $K_R\cdot[\omega]<0$};
  \node[draw=tgCoral,fill=tgCoral!10,rounded corners=4pt,
     minimum width=3.0cm,minimum height=1.1cm,align=center] (c)
     at (0,0) {\small konfiguracja $C_7$\\[-1pt]
             \scriptsize ujemne parowanie};
  \node[draw=tgTeal,fill=tgTeal!10,rounded corners=4pt,
     minimum width=3.1cm,minimum height=1.1cm,align=center] (x)
     at (4.0,0) {\small $X_7$ po kontrakcji\\[-1pt]
             \scriptsize $K_{X_7}\cdot[\omega_7]>0$};
  \draw[-{Stealth[length=2.6mm]},tgPurple,line width=1.7pt] (r)--(c);
  \draw[-{Stealth[length=2.6mm]},tgPurple,line width=1.7pt] (c)--(x);
  \node[text=tgNavy,font=\footnotesize] at (0,-1.05)
      {Wymierna kontrakcja zmienia znak istotnego parowania.};
\end{tikzpicture}
\end{center}

\section{Komory przy \texorpdfstring{$b_2^+=1$}{b2+=1} i małe perturbacje}
\label{sec:tg-park14-chamber}

Przypomnijmy, że dla struktury $\mathrm{Spin}^c$ o klasie
$c:=c_1(L)$ i wymiarze oczekiwanym $d=0$ niezmiennik
Seiberga--Wittena jest liczbą otrzymaną ze zorientowanego
zliczenia rozwiązań dla danych generycznych. Jeśli $b_2^+=1$,
przestrzeń metryk i perturbacji ma ścianę z rozwiązaniami
redukowalnymi. Rozdziela ona dwa rodzaje komór.
Do porównania z metryką dodatniej krzywizny skalarnej
wybieramy perturbacje dostatecznie małe, osobno wzdłuż
każdej zwartej ścieżki metryk; uzyskaną wartość oznaczamy
$SW^{\circ}_X(\mathfrak s)$.

\begin{lemat}[Brak ściany dla małych perturbacji]
\label{lem:tg-park14-no-wall}
Niech $M$ będzie zamkniętą, spójną, orientowaną
czterorozmaitością z $b_2^+(M)=1$ i ustaloną składową
stożka dodatniego. Jeżeli $c\in H^2(M;\mathbb R)$
spełnia $c^2>0$, to jego rzut harmoniczny samodualny
$c_g^+$ jest niezerowy dla każdej metryki $g$.
Co więcej, znak $c\cdot H_g$ dla generatora $H_g$
ustalonej dodatniej półprostej $H_g^2>0$ jest niezależny
od metryki. Jeśli $c$ jest rzeczywistym obrazem klasy wyznacznikowej
$c_1(L)$ ustalonej struktury $\mathrm{Spin}^c$, to
dla perturbacji spełniających
$\|\eta\|_{L^2(g)}<2\pi\sqrt{c^2}$ nie występują
rozwiązania redukowalne. Wszystkie takie dane
należą do tej samej komory.
\end{lemat}

\begin{proof}
Dekompozycja Hodge'a dla metryki $g$ daje ortogonalny
rozkład $c=c_g^++c_g^-$ na dodatnią i ujemną podprzestrzeń
formy przecięcia. Gdyby $c_g^+=0$, mielibyśmy
$c^2=(c_g^-)^2\le0$, sprzecznie z założeniem.
W wymiarze Lorentzowskim dwa wektory o dodatnim kwadracie,
należące do ustalonych składowych stożka, nie są ortogonalne.
Dokładniej, normalizujemy $H_g^2=1$ i zapisujemy
$c=(c\cdot H_g)H_g+c_g^-$. Wtedy
$(c\cdot H_g)^2=c^2-(c_g^-)^2\ge c^2>0$.
Znak iloczynu $c\cdot H_g$ jest stały na wybranej
spójnej składowej stożka dodatniego: jest to znak
ciągłej, nigdzie niezerowej funkcji na tej składowej.
Jeżeli teraz $c$ jest obrazem rzeczywistym $c_1(L)$,
to w konwencji stwierdzenia \ref{prop:tg-sw-wall}
ściana odpowiada zerowaniu funkcji
\[
\mathfrak w_{\mathfrak s}(g,\eta)
=-2\pi c\cdot H_g-\int_M\eta\wedge\omega_g,
\qquad H_g=[\omega_g],\quad \|\omega_g\|_{L^2}=1.
\]
Z nierówności Cauchy'ego--Schwarza wynika
$|\int_M\eta\wedge\omega_g|\le\|\eta\|_{L^2}$.
Wystarczy zatem $\|\eta\|_{L^2}<2\pi\sqrt{c^2}$,
aby znak funkcji ściany nie zmienił się względem
zerowej perturbacji.
Pozostaje wykazać, że wszystkie te dane rzeczywiście
należą do jednej składowej dopełnienia ściany.
Dla ustalonej metryki $g$ droga
$t\mapsto(g,(1-t)\eta)$, $0\le t\le1$,
pozostaje w tej samej otwartej kuli perturbacji.
Dwie dowolne metryki $g_0,g_1$ łączymy drogą
$g_t=(1-t)g_0+tg_1$, która składa się z metryk
riemannowskich, ponieważ wypukła kombinacja
dodatnio określonych form jest dodatnio określona.
Wzdłuż tej drogi bierzemy perturbację identycznie
zerową (samodualną dla każdej metryki).
Z wykazanego oszacowania $\mathfrak w_{\mathfrak s}$
nigdzie na niej nie znika.
Łącząc trzy drogi, otrzymujemy drogę bez ściany
między dowolnymi danymi spełniającymi założenia.
Nie wymagamy regularności rozwiązań wzdłuż tej drogi:
do konstrukcji kobordyzmu stosuje się następnie
generyczną perturbację drogi z ustalonymi końcami.
\end{proof}

\begin{wniosek}[Niezmiennik z małą perturbacją]
\label{cor:tg-park14-small-invariant}
Przy założeniach poprzedniego lematu niech $M$ będzie
po prostu spójna, a $\mathfrak s$ strukturą $\mathrm{Spin}^c$
o wiązce wyznacznikowej $L$, $d(\mathfrak s)=0$
i $c=c_1(L)\otimes1$, $c^2>0$. Wówczas przy ustalonej orientacji
homologicznej niezmiennik $SW_M^{\circ}(\mathfrak s)$
można określić niezależnie od wyboru metryki, używając
małej generycznej perturbacji. Jeżeli na $M$ istnieje
metryka dodatniej krzywizny skalarnej, to
$SW_M^{\circ}(\mathfrak s)=0$.
\end{wniosek}

\begin{proof}
Z twierdzenia \ref{thm:tg-sw12-regular} istnieją
generyczne gładkie perturbacje dowolnie bliskie zeru.
Na mocy lematu \ref{lem:tg-park14-no-wall}
dla każdej metryki wszystkie perturbacje z podanej
tam kuli $L^2$ leżą w jednej i tej samej komorze.
Stwierdzenie \ref{prop:tg-sw12-count} zapewnia
niezależność zorientowanego zliczenia od wyboru
generycznych danych wewnątrz tej komory, dzięki
wersji argumentu kobordyzmowego dla jednej komory,
uzasadnionej po stwierdzeniu \ref{prop:tg-sw-wall}.
Warunek $b_2^+>1$ ze stwierdzenia
\ref{prop:tg-sw-cobordism} nie jest tutaj spełniony;
ścieżkę bez rozwiązań redukowalnych zapewnia
poprzedni lemat. Definiuje to
$SW_M^\circ(\mathfrak s)$.

Niech teraz $g$ będzie metryką o dodatniej krzywiźnie
skalarnej i $m:=\min_M s_g>0$ (korzystamy ze zwartości).
W konwencji twierdzenia \ref{thm:tg-sw12-spinor-bound}
mamy dla każdego rozwiązania, po scałkowaniu
tożsamości Weitzenböcka,
\[
0=\|\nabla_A\psi\|_{L^2}^2
+\int_M\left(\frac{s_g}{4}|\psi|^2
+\frac14|\psi|^4
+\frac12\langle\rho(i\eta)\psi,\psi\rangle\right)
\,\mathrm d\mu_g(x).
\]
Całka z laplasjanu znika z twierdzenia Stokesa.
Ponieważ $\|\rho(i\eta)\|_{\mathrm{op}}\le\sqrt2|\eta|$,
prawa strona jest nie mniejsza niż
\[
\|\nabla_A\psi\|_{L^2}^2+
\left(\frac m4-\frac{\sqrt2}{2}\|\eta\|_{L^\infty}\right)
\|\psi\|_{L^2}^2+\frac14\|\psi\|_{L^4}^4.
\]
Jeżeli $\|\eta\|_{L^\infty}<m/(2\sqrt2)$,
współczynnik przy normie $L^2$ jest dodatni,
więc $\psi=0$. Takie rozwiązanie byłoby
redukowalne. Wybieramy równocześnie
$\|\eta\|_{L^2}<2\pi\sqrt{c^2}$, co wyklucza
je przez poprzedni lemat. Oba warunki określają
otoczenie zera w topologii gładkiej, zawierające
perturbacje generyczne. Dla każdej takiej perturbacji
przestrzeń moduli jest pusta, a zliczenie wynosi zero.
\end{proof}

\begin{uwaga}[Dlaczego nie wystarczy po prostu porównać komór?]
Dla $b_2^+=1$ symplektyczna komora Taubesa nie musi
pokrywać się z komorą, w której pewna metryka dodatniej
krzywizny skalarnej daje pustą przestrzeń moduli.
Porównanie niezerowego zliczenia w jednej komorze
z zerowym zliczeniem w drugiej nie prowadzi do sprzeczności.
W zastosowaniu do Parka potrzebne są oba warunki:
$c^2=K_{X_7}^2=2>0$ zapewnia stałość komory małych
perturbacji, a $c\cdot[\omega_7]<0$ identyfikuje
ją z komorą symplektyczną z następnego twierdzenia.
\end{uwaga}

\section{Twierdzenie Taubesa--Li--Liu i brak dyfeomorfizmu}
\label{sec:tg-park14-exotic}

\begin{twierdzenie}[Taubes--Li--Liu; wynik zewnętrzny]
\label{thm:tg-park14-taubes}
Niech $(M,\omega)$ będzie zamkniętą, spójną,
symplektyczną czterorozmaitością z $b_2^+(M)=1$,
a $K_M=-c_1(TM,J)$ jej klasą kanoniczną.
Dla kanonicznej struktury $\mathrm{Spin}^c$
wyznaczonej przez zgodną strukturę prawie zespoloną,
o wiązce wyznacznikowej $L$ i $c_1(L)=-K_M$,
przy wyborze składowej stożka dodatniego zawierającej
$[\omega]$ oraz ustalonej orientacji homologicznej, niezmiennik w odpowiedniej
komorze symplektycznej spełnia
\[
SW_M^{\mathrm{symp}}(-K_M)=\pm1.
\]
W oznaczeniach Parka (twierdzenie 2.2, s.~3
arXiv:math/0311395v3) jest to $SW_M^-(-K_M)=\pm1$,
gdzie znak komory określa
$2\pi c_1(L)\cdot H_g+\int_M\eta\wedge\omega_g$.
W naszej konwencji z \ref{prop:tg-sw-wall} jest to
$-\mathfrak w_{\mathfrak s}(g,\eta)$:
komora minus Parka odpowiada więc
$\mathfrak w_{\mathfrak s}>0$. Utożsamiamy
oznaczenia perturbacji przez ich występowanie
w równaniu krzywizny, a nie przez samą nazwę komory.
Twierdzenie to przyjmujemy z teorii cechowania;
jego dowód korzysta z odpowiednio przeskalowanych
równań Seiberga--Wittena i analizy granicznej Taubesa.
\end{twierdzenie}

\begin{twierdzenie}[Egzotyczność rozmaitości Parka]
\label{thm:tg-park14-exotic}
Rozmaitość $X_7$ z poprzedniego rozdziału jest
homeomorficzna, lecz nie jest dyfeomorficzna z
$\mathbb{CP}^{2}\#7\overline{\mathbb{CP}}^{\,2}$.
Nie istnieje też na $X_7$ żadna metryka riemannowska
o wszędzie dodatniej krzywiźnie skalarnej.
\end{twierdzenie}

\begin{proof}
Homeomorfizm został już wykazany w twierdzeniu
\ref{thm:tg-park-homeo}. Skupmy się na strukturze gładkiej.
Symplektyczność $X_7$ wynika z geometrycznej konstrukcji
Parka i twierdzenia Symington. Dla klasy kanonicznej $K_7$
wykazaliśmy $K_7\cdot[\omega_7]>0$ w twierdzeniu
\ref{thm:tg-park14-positive}. Zatem
$(-K_7)\cdot[\omega_7]<0$.
Wybieramy składową stożka dodatniego zawierającą
$[\omega_7]$, którego kwadrat
$\int_{X_7}\omega_7\wedge\omega_7$ jest dodatni.
Dla metryki zgodnej z formą symplektyczną
forma $\omega_7$ jest samodualna; ponieważ
jest domknięta, jest również harmoniczna.
Zatem dla tej metryki
$H_g=[\omega_7]/\sqrt{[\omega_7]^2}$.

Ponadto $K_7^2=2e(X_7)+3\sigma(X_7)=2$
na mocy wzoru dla zgodnej struktury prawie zespolonej
i twierdzenia \ref{thm:tg-park-numbers}.
Wymiar oczekiwany moduli dla kanonicznej struktury
$\mathrm{Spin}^c$ o klasie $c_1(L)=-K_7$ wynosi
\[
d(\mathfrak s)=\frac{(-K_7)^2-2e(X_7)-3\sigma(X_7)}4=0.
\]
Ponieważ $(-K_7)^2>0$, lemat
\ref{lem:tg-park14-no-wall} pokazuje, że znak
parowania $(-K_7)\cdot H_g$ pozostaje ujemny dla
każdej metryki i ustalonej składowej stożka dodatniego.
Wobec tego dla małych perturbacji zachodzi
$\mathfrak w_{\mathfrak s}>0$, czyli jest to komora
minus w oznaczeniach Parka. Jest ona właśnie
komorą symplektyczną z twierdzenia
\ref{thm:tg-park14-taubes}. Stąd
\[
SW_{X_7}^{\circ}(-K_7)
 =SW_{X_7}^{\mathrm{symp}}(-K_7)=\pm1.
\]

Gdyby $X_7$ miała metrykę dodatniej krzywizny skalarnej,
wniosek \ref{cor:tg-park14-small-invariant} dałby
$SW_{X_7}^{\circ}(-K_7)=0$, sprzecznie z powyższą równością.
Zatem taka metryka nie istnieje.

Z kolei metryka Fubiniego--Study'ego na
$\mathbb{CP}^{2}$ ma dodatnią krzywiznę skalarną;
odwrócenie orientacji nie zmienia tej własności,
a twierdzenie o zachowywaniu dodatniej krzywizny
skalarnej przy sumie spójnej w wymiarze co najmniej trzy
zapewnia taką metrykę na
$\mathbb{CP}^{2}\#7\overline{\mathbb{CP}}^{\,2}$.
Korzystamy tu z twierdzenia Gromova--Lawsona
o chirurgii zachowującej dodatnią krzywiznę skalarną
(w przypadku sumy spójnej kowymiar chirurgii wynosi
wymiar rozmaitości, a więc tutaj cztery).
Wynik ten przyjmujemy z geometrii riemannowskiej.
Gdyby istniał dyfeomorfizm $F$ z $X_7$ na standardową
rozmaitość z metryką $g_0$ o dodatniej krzywiźnie
skalarnej, to metryka $F^*g_0$ na $X_7$
miałaby krzywiznę skalarną $s_{g_0}\circ F>0$.
Jest to sprzeczne z wykazanym brakiem takiej metryki,
niezależnie od tego, czy $F$ zachowuje orientację.
\end{proof}

\begin{center}
\begin{tikzpicture}[>=Stealth,line cap=round,line join=round]
  \node[draw=tgTeal,fill=tgTeal!10,rounded corners=5pt,
       text width=4.15cm,align=center,inner sep=6pt] (a)
       at (-3.45,0) {\small $X_7$ — konstrukcja Parka\\
          \footnotesize $K_7\cdot[\omega_7]>0$, $(-K_7)^2=2$\\
          \footnotesize $SW^{\circ}(-K_7)=\pm1$};
  \node[draw=tgBlue,fill=tgBlue!10,rounded corners=5pt,
       text width=4.15cm,align=center,inner sep=6pt] (b)
       at (3.45,0) {\small $\mathbb{CP}^{2}\#7\overline{\mathbb{CP}}^{\,2}$\\
          \footnotesize dodatnia krzywizna skalarna\\
          \footnotesize $SW^{\circ}(c)=0$ dla $c^2>0$, $d=0$};
  \node[text=tgCoral,font=\scriptsize,align=center] at (0,.35)
        {nie są\\dyfeomorficzne};
  \node[text=tgPurple,font=\scriptsize,align=center] at (0,-.42)
        {są\\homeomorficzne};
\end{tikzpicture}
\end{center}

\begin{uwaga}[Znaczenie wyniku]
Przykład Parka pokazuje, że forma przecięcia i grupa
fundamentalna, nawet w połączeniu z charakterystyką Eulera
i sygnaturą, nie klasyfikują gładkich czterorozmaitości.
Wymierna kontrakcja zmienia gładką strukturę w sposób,
który nie jest wykrywany przez klasyfikację Freedmana.
Rozstrzygające są dwa dodatkowe fakty: dodatniość
$K_7\cdot[\omega_7]$ oraz niezerowość odpowiedniego
niezmiennika Seiberga--Wittena.
\end{uwaga}

\section{Dokładny status użytych twierdzeń i dalszy krok}
\label{sec:tg-park14-status}

Dowód egzotyczności jest \emph{kompletny warunkowo względem
jawnie wymienionych wyników zewnętrznych}:
(1) istnienia symplektycznej konfiguracji Parka i prostej
spójności jej dopełnienia, (2) twierdzenia Symington,
(3) wyboru klasy symplektycznej z teorii Li--Liu,
(4) twierdzenia Taubesa--Li--Liu i zwartości moduli,
(5) twierdzenia o dodatniej krzywiźnie skalarnej
na sumach spójnych.
Elementarne rachunki form przecięcia, rozkładu klas,
wyznacznika, nierówności i porównania komór zostały
przedstawione szczegółowo. Osobnego wykładu wymagałby
pełny dowód analityczny twierdzenia Taubesa.

W następnej części skryptu będziemy mogli wykorzystać
podobną strategię do operacji chirurgii wzdłuż węzła
(\emph{knot surgery}), gdzie wielomian Alexandera
wchodzi bezpośrednio do formuł na niezmienniki
Seiberga--Wittena. W ten sposób teoria grafów,
teoria węzłów i topologia czterowymiarowa tworzą
jeden ciąg argumentów, a nie trzy odrębne tematy.


% ================================================================
% ITERACJA 15 — rozdzial wstawiany po egzotycznosci Parka
% Konwencje: długosc preferowana lambda_K; zmienna u = exp(2[T]);
% twierdzenia Fintushela--Sterna cytowane bez pelnego dowodu.
% ================================================================
